\section{Sparsification for vertex-induced subgraphs}
\label{sec:spars}
In order to give a space efficient implementation of \cref{alg:metaed}, we will take a subroutine computing an approximate most-balanced sparse cut of a graph, and run it as a black box on a vertex-induced subgraph of a sparsifier. From such sparsifier, we demand that the local cuts are approximated to within an additive error proportional the corresponding global cut. Formally, we rely on the following cut preserving property for vertex-induced subgraphs.
\begin{definition}[Cluster sparsifier]
	\label{def:spars}
	Let $G=(V,E)$, let $\delta \in (0,1)$, let $U \subseteq V$ be a cluster, and let $H=(V,E',w)$ be a weighted subgraph of $G$ on the vertex set $V$. Then, $H$ is a cluster sparsifier for $U$ in $G$ with error~$\delta$, denoted by $H[U] \approx_\delta G[U]$, if one has
	\begin{equation*}
		\forall S \subseteq U,\,\uncut{S}{U} - \delta \cdot \uncutg{S} \le \wcut{S}{U}  \le \uncut{S}{U} + \delta \cdot \uncutg{S} \quad \text{ and }  \quad \forall S \subseteq V,\, (1-\delta)\cdot \uncutg{S}\le \wcutg{S} \le (1+\delta)\cdot \uncutg{S}\, .
	\end{equation*}
\end{definition}
\noindent
In \cref{def:spars}, the second property is the same as the first one when $U=V$. Therefore, if one is able to sample a sparsifier where the first property holds with high probability for any cluster $U\subseteq V$, then we get a sparsifier that with high probability embodies \cref{def:spars}. Perhaps surprisingly, we show that classical constructions of cut sparsifiers give the first property of \cref{def:spars}. Specifically, we use the approach of Fung et al.~\cite{fungspars}, based on sampling edges with probability proportional to their edge connectivity, together with its dynamic stream implementation by Ahn, Guha, and McGregor~\cite{sparsdynamicstreams}. Hence, we obtain the following sparsification lemma.

\begin{restatable}[Cluster sparsifiers in dynamic streams]{lemma}{maintechlemadd}
	\label{lem:distrcutstreamadd}
	Let $G=(V,E)$ be an $n$-vertex graph and let $\delta \in (0,1)$. Then, there is a distribution $\mathcal{D}_\delta$ such that, for any cluster $U \subseteq V$ and a sample $H=(V,E',w)  \sim \mathcal{D}_\delta$ one has $H[U] \approx_\delta G[U]$ with high probability, and $|E'| = \otil(n/\delta^2)$ edges. When $G$ is given in a dynamic stream, there is an algorithm that maintains a linear sketch of $G$ in $\otil(n/\delta^2)$ space and decodes it to output a weighted subgraph $H=(V,E',w)$  in $\otil(n/\delta^2)$ space and $\poly(n)$ time such that $H \sim  \mathcal{D}_\delta$.
\end{restatable}
\noindent
The distribution \smash{$\mathcal{D}_\delta$} we use is a cosmetic modification of the algorithm of Ahn, Guha, and McGregor~\cite{sparsdynamicstreams}, hereafter referred to as the AGM algorithm. It provides a dynamic stream implementation of the sampling scheme of Fung et al.~\cite{fungspars} for constructing cut sparsifiers. Even though we deviate only slightly from the original analysis of Ahn, Guha, and McGregor and Fung et al., we give a complete proof for the sake completeness. In \cref{subsec:primer}, we give a few preliminaries on sparsification and dynamic streams that are needed to present and discuss the algorithm. In \cref{subsec:properties}, we outline the algorithm and prove the technical lemmas that characterise its correctness. In \cref{subsec:cutcounting}, we use standard cut counting arguments to combine such technical lemmas into proving \cref{lem:distrcutstreamadd}.

\subsection{Sparsification primer}
\label{subsec:primer}
We give a few preliminaries needed to present the AGM algorithm. We start with the following result: for a graph given in a dynamic stream, we can output a subgraph that preserves the distinction between $k$-edge connected and not $k$-edge connected cuts.

\begin{theorem}[Connectivity witness \cite{linmeas}]
	\label{th:connwit}
	Let $G=(V,E)$ be a graph given in a dynamic stream, and let $k \ge 1$ be an integer parameter. Then there is an algorithm $\textsc{ConnWit}(k)$ that maintains a linear sketch of $G$ in $\otil(kn)$ space. With high probability, the algorithm decodes the sketch to output a subgraph $G'=(V,E')$ in $\otil(kn)$ space and $\poly(n)$ time such that:
	\begin{enumerate}
		\item for all $\emptyset \neq S \subsetneq V$, if $|E(S,V\setminus S)| < k$ then $E(S, V \setminus S) \subseteq E'$;
		\item for all $\emptyset \neq S \subsetneq V$, if $|E(S,V\setminus S)| \ge k$ then $|E(S, V \setminus S) \cap E'|\ge k$;
		\item $|E'|=O(kn)$.
	\end{enumerate}
\end{theorem}

\begin{proof}
	The algorithm works by maintaining $k$ independent linear sketches for dynamic spanning forest~\cite{linmeas}, each taking $\otil(n)$ space to maintain and $\otil(n)$ space and $\poly(n)$ time to decode. We use the first sketch to recover a spanning forest of $G$. Then, we feed edge removals for every edge in the spanning forest to the second to $k$-th sketch. We proceed in this fashion until we have recovered $k$ edge-disjoint maximal spanning forest $F_1, \dots, F_k \subseteq E$ of $G$. Then we let $E'$ be the union of these spanning forests. This takes $k \cdot \otil(n)$ space and one has $|E'| = O(kn)$.
	
	Now consider a cut $\emptyset \neq S \subsetneq V$. Clearly $E'(S,V \setminus S) \subseteq E(S,V \setminus S)$. Now suppose $E(S,V \setminus S) < k$. For the sake of a contradiction, let us say $E(S,V \setminus S) \setminus E' \neq \emptyset$. By disjointness of the $F_i$'s, there must be a $j \in [k]$ such that the residual graph $G \setminus (\cup_{i=1}^{j-1} F_i)$ still contains an edge from $E(S,V \setminus S)$. Hence, there are two components of the spanning forest $F_j$ which could be joint by this edge, which contradicts its maximality.
	
	Consider the case of $E(S,V \setminus S) \ge k$, and suppose for the sake of a contradiction that $E'(S,V \setminus S) < k$. By disjointness of the $F_i$'s there must be a forest $F_j$ such that $F_j \cap E(S,V \setminus S) = \emptyset$ (as otherwise $E'(S,V \setminus S) \ge k$). As before, this contradicts the maximality of $F_j$.
\end{proof}
\noindent
Hereafter let $C>0$ be a large enough constant, and define for every $e=\{u,v\} \in \binom{V}{2}$ the ideal sampling probabilities
\begin{equation*}
	p_e = \min \left\{1, \frac{1}{\lambda_e}\cdot\frac{C \log^3 n}{\delta^2} \right\} \, ,
\end{equation*}
where $\lambda_e$ is the edge connectivity of $e$, i.e. the minimum number of edges crossing a $uv$-cut in $G$. The original AGM algorithm estimates edge connectivities (and hence sampling probabilities) via the sparsifier that is being constructed, in some sense. Here, we simplify the analysis and instead use a separate sparsifier for that purpose. For instance, one can use a spectral sparsifier, defined as follows.
\begin{definition}[Spectral sparsifier]
	Let $G=(V,E)$ and let $\xi \in (0,1)$. A $(1\pm \xi)$-spectral sparsifier of $G$ is a weighted subgraph $\tilde{G} = (V,\tilde{E},\tilde{w})$ such that
	\begin{equation*}
		\forall \vt{x} \in \mathbb{R}^V , \quad (1-\xi)\vt{x}^\top L \vt{x} \le \vt{x}^\top \tilde{L} \vt{x} \le (1+\xi)\vt{x}^\top L \vt{x} \, ,
	\end{equation*}
	where $L$ and  $\tilde{L}$ are the Laplacian matrices of $G$ and $\tilde{G}$, respectively.
	The Laplacian matrix of a graph is $L=D-A$, where $D$ denotes the degree diagonal matrix of the graph and $A$ denotes its (weighted) adjacency matrix
\end{definition}\noindent
It is known how to construct such sparsifiers in dynamic streams.
\begin{theorem}[\cite{spectralspars}]
	\label{th:spectral}
	Let $G=(V,E)$ be a graph given in a dynamic stream, and let $\xi \in (0,1)$. Then, there is an algorithm $\textsc{Spectral}(\xi)$ that maintains a linear sketch of $G$ in $\otil(n/\xi^2)$ space. With high probability, the algorithm decodes the sketch to output in $\otil(n/\epsilon^2)$ time and space a weighted subgraph  $\tilde{G} = (V,\tilde{E},\tilde{w})$ with $|\tilde{E}| = \otil(n/\epsilon^2)$ that is a $(1\pm \xi)$-spectral sparsifier of $G$.
\end{theorem}

\subsection{Sampling algorithm for dynamic streams}
\label{subsec:properties}
With all the preliminaries in place, we can describe (our version of) the AGM algorithm. The high level approach consists in sampling edges at geometric rates as they arrive in the stream, and then use the $i$-th connectivity witness to recover the edges $e$ that are sampled at rate $p_e \approx 2^i$. Algorithm~\ref{alg:dynsampling} outlines this process, where we denote by $\tilde{\lambda}_e$ the edge connectivity of \smash{$e \in \binom{V}{2}$} in \smash{$\tilde{G}$}. As one can see from the description of the algorithm, we actually let $w$ be a function $E \rightarrow \mathbb{R}_{\ge 0}$ where $w(e)=0$ for all $e \in E \setminus E'$. This facilitates the discussion.

\begin{algorithm}
	\caption{\textsc{agm}: sparsifier construction in dynamic streams}
	\label{alg:dynsampling}
	
	\begin{algorithmic}[1]
		\LeftComment{$\delta \in (0,1)$ is the error parameter}
		\LeftComment{$C >0$ is a large enough constant}
		\LeftComment{$k = 16 C \cdot \delta^{-2}\cdot \log^3 n$}
		\Statex
		\Procedure{\textsc{PreProcessing}}{}
		\State $\textsc{Spectral}\gets $ instance of $\textsc{Spectral}(1/2)$ from \cref{th:spectral}
		\For{$i=1,\dots,{2 \log n}$}
		\State $h_i \sim \unif\left(\{0,1\}^{\binom{V}{2}}\right)$ \Comment{sample $2 \log n$ independent uniform hash functions}
		\State $\textsc{ConnWit}_i(k) \gets $ $i$-th independent instance of $\textsc{ConnWit}(k)$ from \cref{th:connwit} \label{lineagm:feedconnwit}
		\EndFor
		\EndProcedure
		\Statex
		\Procedure{\textsc{OnUpdate}$(e$)}{} \Comment{$e \in \binom{V}{2}$}
		\State feed the update for $e$ to $\textsc{Spectral}$
		\For{$i = 0, \dots, 2\log n$}
		\If{$\prod_{j = 1}^{i} h_j(e)=1$}
		\State feed the update for $e$ to $\textsc{ConnWit}_i(k)$
		\EndIf
		\EndFor
		\EndProcedure
		\Statex
		\Procedure{\textsc{PostProcessing}}{}
		\State $\tilde{G}=(V,\tilde{E},\tilde{w}) \gets $ result of $\textsc{Spectral}$
		\For{$i=0,\dots,{2 \log n}$}
		\State $G_i= (V,E_i) \gets (V,\{e \in E: \prod_{j = 1}^{i} h_j(e)=1\})$ \Comment{used for convenience of analysis only}
		\State $G_i'=(V,E'_i) \gets $ result of $\textsc{ConnWit}_i(k)$ \Comment{used by the algorithm}
		\EndFor
		\For{$e \in \binom{V}{2}$}
		\State $\tilde{p}_e \gets \min \left\{1, \frac{2}{\tilde{\lambda}_e}\cdot\frac{C \log^3 n}{\delta^2} \right\} $
		\State $j_e \gets  \left\lfloor \log \frac{1}{\min\{1, \tilde{p}_e\}}\right\rfloor $.
		\If{$e \in E'_{j_e}$}
		\State $w(e) \gets 2^{j_e}$
		\Else
		\State $w(e)\gets 0$
		\EndIf
		\EndFor
		\State $E' \gets \supp(w)$
		\State \Return $H=(V,E',w)$.
		\EndProcedure
	\end{algorithmic}
\end{algorithm}

We remark that Algorithm~\ref{alg:dynsampling} assumes access to $\poly(n)$ bits of randomness. It is known that one can lift this assumption by using Nisan's pseudorandom generator~\cite{nisansprg,lpindyk} with only a $\pylog(n)$ space blow-up factor (and a $\poly(n)$ time blow-up). We remark that we can do this because the random bits are only used during the stream processing, and we only compute linear sketches in this phase, which in particular are oblivious to the order of updates in the stream. Then we can already conclude a bound on the space complexity, as well as on the size of the output graph.
\begin{lemma}
	\label{lem:spacesampling}
	Algorithm~\ref{alg:dynsampling} takes $\otil(n/\delta^2)$ bits of space and $H$ has at most $\otil(n/\delta^2)$ edges.
\end{lemma}
\begin{proof}
	The edges of $H$ can only come from the output of one of the instances $\textsc{ConnWit}$. In other words, $E' \subseteq \cup_i E_i'$. From \cref{th:connwit} we also know that each instance of $\textsc{ConnWit}$ takes $\otil(kn)$ bits of space and that each $E_i'$ contains at most $\otil(kn)$ edges. Hence, $|E'| = \otil(kn)=\otil(n/\delta^2)$ and Algorithm~\ref{alg:dynsampling} takes \smash{$\otil(kn) = \otil(n/\delta^2)$} bits of space.
\end{proof}

\noindent
We continue our analysis by bounding the probability that a certain bad event happens, which corresponds to the sampling or sketching primitives failing to preserve a needed value. To be more specific about this, first define for every $i =0, \dots, 2\log n$ the event $\mathcal{B}_i$ that instance $i$ of Theorem~\ref{th:connwit} $\textsc{ConnWit}_i(k)$ fails to correctly construct the graph $G_i'$. In addition, for every $e \in \binom{V}{2}$, let \smash{$\mathcal{B}^{(1)}_e$} be the event that \smash{$\tilde{\lambda}_e < 1/2 \cdot \lambda_e$ or $\tilde{\lambda}_e > 3/2 \cdot \lambda_e$}. Finally, define for every $e \in E$ the quantity
\begin{equation*}
	\tau_e = \left\lfloor \log \frac{1}{\min\{1, p_e\}}\right\rfloor \, ,
\end{equation*} and let \smash{$\mathcal{B}^{(2)}_e$} be the event that there exists $j  \in \mathbb{N}$ with $\tau_e-2 \le j \le \tau_e$ such that $\lambda_e(G_{j}) \ge k$, where $k$ is defined in \cref{alg:dynsampling} and $\lambda_e(G_i)$ is the edge connectivity of $e$ in $G_i$ for $i=0,\dots,2\log n$. With this notation, the bad event we want to avert is
\begin{equation*}
	\mathcal{B} = \Big(\vee_{i=0}^{2 \log n} \mathcal{B}_i\Big) \vee \Big(\vee_{e \in \binom{V}{2}} \mathcal{B}^{(1)}_e\Big) \vee \Big(\vee_{e \in E} \mathcal{B}^{(2)}_e\Big) \, .
\end{equation*}
Intuitively, this is a bad event because the algorithm cannot trust the connectivity witnesses nor the edge connectivities computed via the spectral sparsifiers. On the contrary, when $\mathcal{B}$ does not happen, we know that all the graphs $G_i'$ are constructed correctly (thanks to $\neg  \mathcal{B}_i$), all the edges have $j_e$ that is not too large (thanks to \smash{$\neg  \mathcal{B}_e^{(1)}$}), and we can recover the sampled edges via the graphs~$G_i'$ (thanks to \smash{$\neg  \mathcal{B}_e^{(2)}$}).

\begin{lemma}
	\label{lem:failsampling}
	When running Algorithm~\ref{alg:dynsampling}, the bad event $\mathcal{B}$ does not happen with high probability.
\end{lemma}
\begin{proof}
	By \cref{th:spectral}, we know that $\vee_{e} \mathcal{B}^{(1)}_e$ does not happen with high probability for any $e \in \binom{V}{2}$. Now let $i=0,\dots, 2 \log n$, and note that the randomness used to define the input $G_i$ to $\textsc{ConnWit}_i(k)$ and the randomness used by $\textsc{ConnWit}_i(k)$ itself are independent. Hence, by \cref{th:connwit} and by union bound over $i$, we get that $\vee_{i}\mathcal{B}_i$ does not happen with high probability.  Lastly, we deal with the events $\mathcal{B}^{(2)}_e$.
	
	For ease of notation, define for any $\emptyset \neq S \subseteq V$ the size of its cut projected onto the edges of $G_i$ as $\partial^i S=|E_i \cap E(S,V \setminus S)|$. Then consider any $e=\{s,t\} \in E$ and any $j  \in \mathbb{N}$ with $\tau_e-2 \le j \le \tau_e$, and let $S_e \subseteq V$ be an $st$-cut in $G$ of minimum size, i.e. $\lambda_e = \uncutg{S_e}$. Observe that by the way the algorithm defines the sub-streams, each $f \in E$ is in $E_i$ with probability $2^{-i}$ for all $i=0, \dots, 2 \log n$, so in expectation the edge connectivity of $e$ has dropped below $4 C \cdot \delta^{-2}\cdot \log^3 n$ in the graph $G_{j}$. More precisely,
	\begin{align*}
		\ex[\partial^{j} S_e] = \uncutg{S_e} \cdot 2^{-j} & \le 4\cdot \uncutg{S_e} \cdot 2^{-\tau_e} \\
		&  \le 8\lambda_e \min\{1,  p_e\} \\
		& \le 8\lambda_e p_e  \\
		& \le 8 \lambda_e \frac{C \cdot \delta^{-2}\cdot \log n}{\lambda_e} = 8 C \cdot \delta^{-2}\cdot \log^3 n \, .
	\end{align*}
	Applying the Chernoff bound (since edges are sampled independently of each other)  we get that
	\begin{align*}
		\Pr\left[\partial^{j} S_e > 12 C \cdot \delta^{-2}\cdot \log^3 n \right] & \le \Pr\left[\partial^{j} S_e > \left(1+\frac{1}{2}\right) \ex[\partial^{j} S_e] \right] \\
		& \le \exp\left(-\frac{\ex[\partial^{j} S_e]}{12}\right) \\
		& \le \exp\left(-\frac{\ex[\partial^{\tau_e} S_e]}{12}\right)\, .
	\end{align*}
	In order to make sure that the exponential above is in fact small enough, we now have two cases to consider. First consider the possibility that $\tau_e= 0$, or equivalently $p_e > \frac{1}{2}$. This means two things: $E_{\tau_e}=E$ (see \cref{alg:dynsampling}), and $\lambda_e < 2 C \cdot \delta^{-2}\cdot \log^3 n$. Then, $\partial^{\tau_e} S_e=\lambda_e \le 12 C \cdot \delta^{-2}\cdot \log^3 n$ with probability $1$, so we need not worry about this case. For the other possibility that $ p_e \le \frac{1}{2}$, one can deduce that $p_e = C \cdot \lambda_{e}^{-1} \delta^{-2}\cdot \log^3 n$ and thus lower bound the expectation of $\partial^{\tau_e} S_e$ as
	\begin{equation*}
		\ex[\partial^{\tau_e} S_e] = \partial S_e \cdot 2^{-\tau_e} \ge \lambda_e p_e = C \cdot \delta^{-2}\cdot \log^3 n \, ,
	\end{equation*}
	Hence, the Chernoff bound above gives $\partial^{\tau_e} S_e \le 12 C \cdot \delta^{-2}\cdot \log^3 n < k$ with high probability. Thus, the edge connectivity of $e$ in $G_{j}$ must be at most $\partial^{j} S_e < k$. Taking a union bound gives that $\vee_{e} \mathcal{B}^{(2)}_e$ does not happen with high probability.
\end{proof}

\noindent
Finally, we show properties of the weight function produced by the algorithm. We first focus on the special case of edges with very low connectivity, i.e. with $p_e=1$, and show that these are indeed copied exactly in $E'$, as one would expect.

\begin{lemma}
	\label{lem:oneprobsampling}
	Let $H=(V,E',w)$ be the output of Algorithm \ref{alg:dynsampling}. If $\mathcal{B}$ does not happen, then for any edge $e \in E$ with $p_e =1$ we have $e \in E'$ and $w(e)=1$ .
\end{lemma}
\begin{proof}
	Assume $\mathcal{B}$ does not happen. This in particular means that $G_0'$ is correctly constructed, so by Theorem~\ref{th:connwit} any cut of size less than $16C \cdot \delta^{-2}\cdot \log^3$ is exactly preserved in $G_0'$. On the other hand, if an edge $e \in E$ has $p_e = 1$ then its edge connectivity $\lambda_e$ is at most $C \cdot \delta^{-2}\cdot \log^3 n$. Since $\tilde{\lambda}_e \le 2 \lambda_e$, it then follows that $\tilde{p}_e = 1 $, so $j_e=0$ and hence the claim.
\end{proof}

\noindent
For edges which we do not know to have $p_e=1$, but for which we at least know a lower bound $p$ on their sampling probability, we wish to show that they are preserved with ``very high'' probability. We do this by using a Chernoff-style bound of Ahn, Guha, and McGregor.

\begin{restatable}{lemma}{probsampling}
	\label{lem:probsampling}
	Let $H=(V,E',w)$ be the output of Algorithm \ref{alg:dynsampling}. Then for any edge set $X \subseteq E$, any $p \le \min_{e \in X} p_e$, and any $x \ge |X|$ one has
	\begin{equation*}
		\Pr\left[\neg \mathcal{B} \text{ and } \left|\sum_{e \in X} w(e) - |X|\right| > \delta \cdot x \right] \le 2 \exp\left(-\frac{1}{48} \cdot \delta^2 \cdot  p x\right).
	\end{equation*}
\end{restatable}
\begin{proof}
	For ease of notation, let $\tilde{\lambda}$ be the vector of edge connectivities $(\tilde{\lambda}_e)_{e \in E}$ from \cref{alg:dynsampling}, and $\lambda$ is the vector of edge connectivities $(\lambda_e)_{e\in E}$ of $G$. We will condition on the values $\hat{\lambda}$ that $\tilde{\lambda}$ can take, and define
	\begin{equation*}
		\hat{p}_e = \min \left\{1, \frac{2}{\hat{\lambda}_e}\cdot\frac{C \log^3 n}{\delta^2} \right\} \quad \text{ and } \quad  \hat{j}_e = \left\lfloor \log \frac{1}{\min\{1, \hat{p}_e\}}\right\rfloor  \, .
	\end{equation*}
	Then,
	\begin{align*}
		& \Pr\left[\neg \mathcal{B} \text{ and } \left|\sum_{e \in X} w(e) - |X|\right| > \delta \cdot x \right] \\
		= & \sum_{\hat{\lambda} \in \mathbb{N}^E} 	\Pr\left[\neg \mathcal{B} \text{ and } \left|\sum_{e \in X} w(e) - |X|\right| > \delta \cdot x  \text{ and } \tilde{\lambda}=\hat{\lambda} \right] \\
		= & \sum_{\hat{\lambda} \in \mathbb{N}^E \,:\, \frac{1}{2} \lambda \le {\hat{\lambda}} \le \frac{3}{2} \lambda	} \Pr\left[\neg \mathcal{B} \text{ and } \left|\sum_{e \in X} w(e) - |X|\right| > \delta \cdot x  \text{ and } \tilde{\lambda}=\hat{\lambda} \right] \, .
	\end{align*}
	By conditioning we get
	\begin{align*}
		& \Pr\left[\neg \mathcal{B} \text{ and } \left|\sum_{e \in X} w(e) - |X|\right| > \delta \cdot x \right] \\
		= & \sum_{\hat{\lambda} \in \mathbb{N}^E \,:\, \frac{1}{2} \lambda \le {\hat{\lambda}} \le \frac{3}{2} \lambda	} \Pr\left[\neg \mathcal{B}\text{ and } \left|\sum_{e \in X} w(e) - |X|\right| > \delta \cdot x  \middle| \tilde{\lambda}=\hat{\lambda} \right] \cdot \Pr\left[\tilde{\lambda}=\hat{\lambda}\right]
	\end{align*}
	We now analyse the first factor in each term of the sum above. We use three facts to do this: (1) whenever we have $\neg \mathcal{B}$, for all $e \in X$ one has that $e \in E_{j_e}'$ if and only if $e \in E_{j_e}$ (this is the case because \smash{$\frac{1}{2} \lambda \le {\tilde{\lambda}} \le \frac{3}{2} \lambda$}, so \smash{$\max\{0, \tau_e-2\} \le j_e \le \tau_e$, and the event $\mathcal{B}^{(2)}_e$} does not happen); (2) recall that, for $i=0,\dots,2\log n$, $e \in E_{i}$ if and only if $\prod_{t=1}^{i}h_t(e)=1$ (see \cref{alg:dynsampling}); (3) also recall that \cref{alg:dynsampling} assigns weight $2^{j_e}$ to an edge recovered from $E'_{j_e}$ (see \cref{alg:dynsampling}).
	Using these three facts, we get
	\begin{equation*}
		 \Pr\left[\neg \mathcal{B} \text{ and } \left|\sum_{e \in X} w(e) - |X|\right| > \delta \cdot x  \middle| \tilde{\lambda}=\hat{\lambda} \right]
		= \Pr\left[\neg \mathcal{B} \text{ and } \left|\sum_{e \in X} 2^{j_e}\prod_{t=1}^{j_e}h_t(e) - |X|\right| > \delta \cdot x  \middle| \tilde{\lambda}=\hat{\lambda} \right] \, .
	\end{equation*}
	Next, by the condition $\tilde{\lambda}=\hat{\lambda}$, we have $\hat{j}_e=j_e$ for all $e \in X$, so we get
	\begin{align*}
	\Pr\left[\neg \mathcal{B} \text{ and } \left|\sum_{e \in X} w(e) - |X|\right| > \delta \cdot x  \middle| \tilde{\lambda}=\hat{\lambda} \right]
		& =  \Pr\left[\neg \mathcal{B} \text{ and } \left|\sum_{e \in X} 2^{\hat{j}_e}\prod_{t=1}^{\hat{j}_e}h_t(e) - |X|\right| > \delta \cdot x  \middle| \tilde{\lambda}=\hat{\lambda} \right] \\
		& \le \Pr\left[\left|\sum_{e \in X} 2^{\hat{j}_e}\prod_{t=1}^{\hat{j}_e}h_t(e) - |X|\right| > \delta \cdot x  \middle| \tilde{\lambda}=\hat{\lambda} \right] \\
		& = \Pr\left[\left|\sum_{e \in X} 2^{\hat{j}_e}\prod_{t=1}^{\hat{j}_e}h_t(e) - |X|\right| > \delta \cdot x \right]  \, .
	\end{align*}
	Now we apply the following form of the Chernoff bound.
	\begin{fact}[\cite{streamexpdec}]
		Let $X_1,\dots,X_N$ be independent random variables distributed in $[0,a]$. Let $X= \sum_{r \in [N]} X_r$ and $\mu = \ex[X]$. Then, for $\zeta \in (0,1)$ and $\alpha \ge 0$ one has
		\begin{equation*}
			\Pr\left[\left| X- \mu \right| > \zeta \mu+\alpha \right] \le 2 \exp\left(-\frac{\zeta \alpha}{3a}\right)
		\end{equation*}
	\end{fact}
	\noindent
	Before applying this bound, we note that the random variables $2^{\hat{j}_e}\prod_{t=1}^{\hat{j}_e}h_t(e)$ are independent across $e \in X$ and have each mean $1$. Moreover, each of this random variables have value at most \smash{$2^{\hat{j}_e} \le 2^{\lfloor \log \frac{1}{\min\{1,p\}}\rfloor}$}, so we set $a = 1/p$. We also set $\zeta = \delta/2$ and $\alpha = \delta\cdot (x-|X|)+\delta|X|/2$. Then we get
	\begin{equation*}
		\Pr\left[\left|\sum_{e \in X} 2^{\hat{j}_e}\prod_{t=1}^{\hat{j}_e}h_t(e) - |X|\right| > \delta \cdot x \right] \le 2 \exp\left(-\frac{\delta^2}{48}px\right) \, .
	\end{equation*}
	This concludes the proof.
\end{proof}

\subsection{Cut counting and proving the sparsification lemma}
\label{subsec:cutcounting}
We hereafter fix a cluster $U\subseteq V$, and we partition the edge set $E$ into $F_i = \{e \in E: \, 2^i \le \lambda_e \le 2^{i+1} -1 \, \}$ for $i =0,\dots,\log n$. Note that these $F_i$'s are pairwise disjoint and their union gives~$E$. We recall that $E(A,B)$ denotes the set of edges in $E$ with one endpoint in $A$ and one in $B$, for $A,B \subseteq V$. Then, we write $\iuncut{S}{U}=|F_i \cap E(S,U \setminus S)|$ to denote the size of the projection of the local cut of $S$ onto the edge set $F_i$, for all $i=0,\dots, \log n$ and $\emptyset \neq S \subsetneq U$. Analogously, we also write $\iwcut{S}{U}=w(F_i \cap E(S,U \setminus S))$ to denote the total weight of the projection of the local cut of $S$ onto~$F_i$, where $w$ is the weight function of $H$.

We show that $H$ preserves the cuts of a cluster $U \subseteq V$ within each $F_i$ individually, with high probability. This follows by two of the technical lemmas we proved in the previous section, namely \cref{lem:oneprobsampling} and \cref{lem:probsampling}.

\begin{lemma} \label{lem:levelbound}
	Let $H=(V,E',w)$ be the output of \cref{alg:dynsampling}, and let $i =0,\dots,\log n$. Then
	\begin{align*}
		\Pr\left[ \neg B\,  \text{ and }\, \exists \, \emptyset \neq S \subsetneq U \, : \,  \left|\iwcut{S}{U} - \iuncut{S}{U}\right| > \delta \cdot \frac{\uncutg{S}}{1+\log n} \right] \le \frac{1}{\poly(n)}
		\, .
	\end{align*}
\end{lemma}
\begin{proof}
We proceed by first showing a probability bound for each cut $S$, and then combine them using a union bound to get the final result. Specifically, for all $j \in [2\log n]$ we define the set
\begin{equation*}
	\mathcal{Y}_{ij}=\{E(S,V \setminus S) : \, \emptyset \neq S \subsetneq U \text{ and } 2^{i+j-1} \le \uncutg{S} \le 2^{i+j}-1 \} \, ,
\end{equation*}
and we will show a probability bound for all $j \in [2\log n]$ and $Y \in \mathcal{Y}_{ij}$. Taking a careful union bound over all such choices will be enough to get the lemma statement. It shall be useful to define $\Pi_U^i(Y)=Y \cap F_i \cap \binom{U}{2}$ to be the projection onto $F_i$ and onto the cluster $U$ of the edges in $Y$, for any $Y \subseteq E$. Moreover, we denote by $\mathcal{F}_{Y}$ the failure event that
\begin{equation*}
	\left|w(\Pi_U^{i}(Y)) - |\Pi_U^{i}(Y)|\right| > \delta \frac{|Y|}{1+\log n} \, ,
\end{equation*}
For ease of writing we shall denote $\Pi_U^i(E(S,V \setminus S))$ and $\mathcal{F}_{E(S,V \setminus S)}$ by $\Pi_U^i(\uncutg{S})$ and $\mathcal{F}_{\uncutg{S}}$ respectively. To appreciate their purpose, we remark that for any $\emptyset \neq S \subsetneq U$,
\begin{align*}
	\mathcal{F}_{\uncutg{S}} \quad \iff \quad \left|\iwcut{S}{U} - \iuncut{S}{U}\right| > \delta \frac{\uncutg{S}}{1+\log n} \, 
\end{align*}
since $\Pi_U^{i}(\uncutg{S})=\Pi_U^{i}(E(S,V \setminus S)) = E(S,V \setminus S) \cap F_i \cap \binom{U}{2} = F_i \cap E(S,U\setminus S)$.

\begin{claim}
	\label{claim:boundprobsingleset}
	Let $j \in [2\log n]$ and $Y \in \mathcal{Y}_{ij}$. Then
	\begin{equation*}
		\Pr[\neg \mathcal{B} \text{ and } \mathcal{F}_Y] \le 2n^{-\frac{C}{768}2^j} \, .
	\end{equation*}
\end{claim}

\begin{proof}
	We start by observing that we can consider only the set $X = \{e \in \Pi_U^i(Y) \, : \, p_e < 1\}$, since edges with $p_e=1$ are always added to $E'$ with weight set to $1$ (when $ \mathcal{B}$ does not occur) by Lemma~\ref{lem:oneprobsampling}, so
	\begin{equation*}
		\left|w(\Pi_U^{i}(Y)) - |\Pi_U^{i}(Y)|\right| = \left|w(X) - |X|\right| \, .
	\end{equation*}
	If $X = \emptyset$, the probability of $\mathcal{F}_Y$ is zero, so let us consider the case that $X \neq \emptyset$, for which we want to upper bound the probability
	\begin{equation*}
		\Pr\left[\neg \mathcal{B} \text{ and } \left|w(X) - |X|\right| > \delta \frac{|Y|}{1+\log n}\right] \, .
	\end{equation*}
	\noindent
	To do so, we use Lemma~\ref{lem:probsampling} and set $p=\min \{1, C \cdot 2^{-(i+1)} \delta^{-2} \log^3 n\}$, which is a lower bound on all the probabilities $\{p_e\}_{e \in X}$. This is the case because every $e \in X$ has $p_e < 1$ and $\lambda_e \le 2^{i+1}$ (since $X \subseteq F_i$). We also note that $|Y| \ge |X|$, thus we can safely apply Lemma~\ref{lem:probsampling} and get
	\begin{align*}
		\Pr\left[\neg \mathcal{B} \text{ and } \left|w(X) - |X|\right| > \delta \frac{|Y|}{1+\log n}\right] & \le 2 \exp\left(-\frac{1}{48}\cdot \delta^2 \cdot p \cdot \frac{|Y|}{(1+\log n)^2} \right) \, .
	\end{align*}
	If we recall $Y \in \mathcal{Y}_{ij}$, and in particular $|Y| \ge 2^{i+j-1}$, we can then conclude the claim by plugging the value of $p$ into the exponent.
\end{proof}
\noindent
Now that we have bounded the probability that a single cut from $\mathcal{Y}_{ij}$ is not well preserved once projected onto $U$ and $F_i$, we show that the union over $j$ of the collections $\mathcal{Y}_{ij}$ is representative of all the local cuts of $U$.
\begin{claim}
	\label{claim:calyisenough}
	Let $\emptyset \neq S \subsetneq U$ such that $\Pi_U^i(\uncutg{S})\neq~\emptyset$. Then there exists $j \in [2\log n]$ so that $ E(S,V \setminus~S)~\in~\mathcal{Y}_{ij}$.
\end{claim}
\begin{proof}
	Since $\Pi_U^i(\uncutg{S}) \neq~\emptyset$, there exists $e=\{u,v\} \in E(S,V\setminus S)$ that is also in $F_i$. In particular this means that $E(S,V\setminus S)$ is a $uv$-cut in $G$, and $e$ has edge connectivity at least $2^i$, thus $\uncutg{S} \ge 2^{i}$. At the same time, $\uncutg{S} \le 2^{i+j}-1$ for some $j \in [2\log n]$ since $\uncutg{S} \le |E| \le n^2-1$.
\end{proof}
\noindent
Now, for all $j \in [2 \log n]$ and all $Z \subseteq F_i$, we define the worst-case edge set $Y^*_j(Z) \subseteq E$ of $Z$ as
\begin{equation*}
	Y^*_j(Z) \in \argmin_{Y \in \mathcal{Y}_{ij}: \Pi_U^i(Y) = Z}  |Y|\, .
\end{equation*}
The reason why we call $Y^*_j(Z)$ worst case for $Z$ is the following. Among all the global cuts in $\mathcal{Y}_{ij}$ whose projection onto $F_i$ and $U$ is $Z$,  the cut \smash{$Y^*_j(Z)$} is the one which allows for the smallest error term, which intuitively makes $\mathcal{F}_Y$ more likely to happen. Formally we have the following claim, which follows from the definition.

\begin{claim}
	\label{claim:worstcaseproj}
	If there exist $j \in [2\log n]$ and $Y \in \mathcal{Y}_{ij}$ such that the event $\mathcal{F}_{Y}$ happens, then there exist $j \in [2\log n]$ and $Z \in \{\Pi_U^i(Y'): Y' \in \mathcal{Y}_{ij}\}$ such that $\mathcal{F}_{Y_j^*(Z)}$ happens.
\end{claim}
\noindent
Claim~\ref{claim:calyisenough} and~\ref{claim:worstcaseproj} together suggest that we can restrain ourselves to union bound over $\{\Pi_U^i(Y): Y \in \mathcal{Y}_{ij}\}$, which is exactly what we are going to do. Before that, we make sure that this set is not too large.
\begin{claim}
	\label{claim:smallunionbound}
	For every $j \in [2\log n]$ we have $|\{\Pi_U^i(Y): Y \in \mathcal{Y}_{ij}\}| \le n^{2 \cdot 2^{j}}$.
\end{claim}
\begin{proof}
	We will employ the following cut counting result, where $E^k$ denotes the set of edges $e \in E$ with edge connectivity $\lambda_e \ge k$, for any integer $k$.
	\begin{theorem}[\cite{fungspars} Theorem 1.6]\label{th:counting}
		For any $k, s \ge 1$ one has
		\begin{equation*}
			|\{E(S,V\setminus S) \cap E^k \, : \, \emptyset \neq S \subsetneq V \text{ and } \uncutg{S} \le s\cdot k\}| \le n^{2s} \, ,
		\end{equation*}
		i.e. the number of distinct subsets of edges of connectivity at least $k$ from cuts of size at most $s \cdot k$ is at most $n^{2s}$.
	\end{theorem}
	\noindent
	Fix $j \in [2\log n]$. We recall that every $e \in F_i$ has $\lambda_e \ge 2^i$, which means $F_i \subseteq E^{2^{i}}$. Also, for any $Y \in \mathcal{Y}_{ij}$ there is by definition a cut $\emptyset \neq S \subsetneq V$ such that $\uncutg{S} \le 2^{i+j}-1$ and $Y=E(S,V\setminus S)$. Hence
	\begin{align*}
		|\{\Pi_U^i(Y) : \, Y \in \mathcal{Y}_{ij}\}|
		& \le |\{E^{2^i} \cap E(S,V\setminus S) : \emptyset \neq S \subsetneq V \text{ and } \uncutg{S} \le 2^{i+j}\}| \, ,
	\end{align*}
	and applying Theorem~\ref{th:counting} with $k \coloneqq 2^i$ and $s \coloneqq 2^j$ we conclude the claim.
\end{proof}
\noindent
Using the above claims we finally have
\begin{align*}
	& \Pr[\neg\mathcal{B} \text{ and } \exists \, \emptyset \neq S \subsetneq U \text{ s.t. } \mathcal{F}_{\uncutg{S}}] \\
	\text{by Claim~\ref{claim:calyisenough}} \quad  \quad \le & \Pr[\exists j \in [2\log n] \text{ s.t. } \exists Y \in \mathcal{Y}_{ij} \text{ s.t. } \neg\mathcal{B} \text{ and } \mathcal{F}_{Y}] \\
	\text{by Claim~\ref{claim:worstcaseproj}} \quad  \quad \le & \Pr[\exists j \in [2\log n] \text{ s.t. } \exists Z \in \{\Pi_U^i(Y): Y \in \mathcal{Y}_{ij}\} \text{ s.t. }  \neg\mathcal{B} \text{ and } \mathcal{F}_{Y_j^*(Z)}] \\
	\text{by Claim~\ref{claim:boundprobsingleset}} \quad  \quad \le & \sum_{j \in [2\log n]} \quad \sum_{ Z \in \{\Pi_U^i(Y) : \, Y \in \mathcal{Y}_{ij}\}} 2n^{-\frac{C}{768}2^j} \\
	\text{by Claim~\ref{claim:smallunionbound}} \quad  \quad \le & \sum_{j \in [2\log n]} n^{2 \cdot 2^{j}}\, 2n^{-\frac{C}{768}2^j} \\
	\le & \frac{1}{\poly(n)}
\end{align*}
\end{proof}
\noindent
Finally, we can conclude the main sparsification lemma \cref{lem:distrcutstreamadd} by a simple combination of \cref{lem:failsampling}, \cref{lem:levelbound}, and \cref{lem:spacesampling}.

\begin{proof}[Proof of \cref{lem:distrcutstreamadd}]
Recall that the $F_i$'s partition and cover $E$ entirely, and that $\neg\mathcal{B}$ happens with high probability by \cref{lem:failsampling}. Then, one can see that conditioning on $\neg \mathcal{B}$, summing over $i=0,\dots, \log n$, and applying \cref{lem:levelbound} gives \cref{lem:distrcutstreamadd}.
\end{proof}